\documentclass{article}
\usepackage[T1]{fontenc}
\usepackage{lmodern}
\usepackage{graphicx}
\usepackage{amsmath,amssymb}
\usepackage[a4paper,margin=2cm]{geometry}
\usepackage[absolute,overlay]{textpos}
\setlength{\TPHorizModule}{1mm}
\setlength{\TPVertModule}{1mm}
\usepackage[indonesian]{babel}
\usepackage{hyperref}
\setlength{\emergencystretch}{2em}
% unit: Halaman 61

\begin{document}

% === Halaman 61 (llm) ===

\begin{itemize}
    \item Dengan menggunakan kunci SCRAM cipherteks berhasil didekripsi menjadi:
    
    \begin{center}
    \begin{tabular}{l}
    THEBE ARWEN TOVER THEMO UNTAI NYEAH \\
    THEDO GWENT ROUND THEHY DRANT THECA \\
    TINTO THEHI GHEST SPOTH ECOUL DFIND
    \end{tabular}
    \end{center}
    
    \item Lakukan post-processing untuk menghasilkan kalimat yang lebih jelas:
    
    \begin{center}
    \begin{tabular}{l}
    THE BEAR WENT OVER THE MOUNTAIN YEAH \\
    THE DOG WENT ROUND THE HYDRANT \\
    THE CAT INTO THE HIGHEST SPOT HE COULD FIND
    \end{tabular}
    \end{center}
\end{itemize}

\end{document}
